\section{Conclusions}
\label{sec:conclusion}

The aim of this study is to determine a useful role for reusable learning in alignment-based conformance checking while keeping the meaning of an explanation explicit. LARA combines a pretrained graph-and-trace scorer with a marking-aware decoder, independent replay, and optional exact comparison. The scorer proposes move preferences; replay establishes executability and an upper bound; successful exact computation establishes the optimum or supplies a replacement.

On the controlled synthetic test set, guidance increases optimal candidates from 445 to 466 out of 512, with its largest contribution on duplicate-prefix choices. All evaluated candidates pass replay, but candidate quality falls on larger stress inputs. Real-log experiments show that the fixed checkpoint can be used with new vocabularies and model structures, while providing no quality improvement over the unguided decoder on those logs. Established approximations generally perform better in the direct comparisons, and the current certified mode still runs exact search for every trace.

The resulting contribution is a transparent, reproducible study of learned candidate scoring within an alignment pipeline whose guarantees come from process semantics. The most promising next steps are stronger decisions among competing move types, broader tests of structural transfer, and using verified incumbents to reduce exact search work. Their value must be demonstrated in candidate quality and total computation, with the same explicit separation between a plausible prediction and a proven alignment property.
